\documentclass[11pt]{article}
\usepackage[a4paper,margin=21mm]{geometry}
\usepackage[no-math]{fontspec}
\setmainfont{Latin Modern Roman}
\newfontfamily\CJKfont{Noto Serif CJK SC}
\XeTeXlinebreaklocale "zh"
\XeTeXlinebreakskip=0pt plus 1pt
\usepackage{amsmath,amssymb,amsthm,amscd,graphicx,color}
\usepackage{xurl}
\usepackage[unicode,hidelinks]{hyperref}
\allowdisplaybreaks
\setlength\emergencystretch{3em}
\numberwithin{equation}{section}
\newtheorem{Theorem}{Theorem}[section]
\newtheorem*{Theorem*}{Theorem}
\newtheorem{Corollary}[Theorem]{Corollary}
\newtheorem{Lemma}[Theorem]{Lemma}
\newtheorem{Proposition}[Theorem]{Proposition}
\newtheorem{Conjecture}[Theorem]{Conjecture}
\theoremstyle{definition}
\newtheorem{Definition}[Theorem]{Definition}
\newtheorem{Note}[Theorem]{Note}
\newtheorem{Example}[Theorem]{Example}
\newtheorem{Remark}[Theorem]{Remark}

\numberwithin{equation}{section}


\renewcommand{\k}{\mathbbm{k}}
\newcommand{\F}{\mathbf{F}}

\newcommand{\sgn}{\mathrm{sign}}
\newcommand{\Irrep}{\mathrm{Irrep}}
\newcommand{\C}{\mathbb{C}}
\newcommand{\Z}{\mathbb{Z}}

\newcommand{\triv}{\mathrm{triv}}
\newcommand{\std}{\mathrm{std}}
\newcommand{\sign}{\mathrm{sign}}
\newcommand{\trix}{\mathrm{triv-sign-triv}}
\newcommand{\six}{\mathrm{sign-triv-sign}}
\newcommand{\siv}{\mathrm{sign-triv}}
\newcommand{\tign}{\mathrm{triv-sign}}
\newcommand{\trtr}{\mathrm{triv-triv}}
\newfontfamily\nativebbk{DejaVu Math TeX Gyre}
\ExplSyntaxOn
\NewDocumentCommand{\mathbbm}{m}{\str_if_eq:nnTF{#1}{k}{\mathord{\text{\mdseries\upshape\nativebbk 𝕜}}}{\PackageError{nativecompat}{Unsupported~mathbbm~argument}{Only~verified~lowercase~k~is~accepted}}}
\ExplSyntaxOff\usepackage{enumitem}
\makeatletter\newlabel{wangv2:sec:4}{{4}{}{Original preprint Section4 preview}{}{}}\makeatother
\begin{document}
\begin{center}\Large Free generators of S3 quasi-invariants in characteristic two\end{center}
\noindent\textbf{Stable ID:} 682\quad\textbf{Author:} Frank Wang; Eric Yee\par
\noindent\textbf{Work:} Hilbert Series of S3-Quasi-Invariant Polynomials in Characteristics2,3\par
\noindent\textbf{Source:} \url{https://sigma-journal.com/2025/057/}\par
\noindent\textbf{Version:} Official SIGMA 21 (2025) journal PDF and native TeX; acquired 2026-09-30; exact bytes pinned by SHA256\par
\noindent\textbf{Rights:} CC BY 4.0. Authors retain copyright. \url{https://creativecommons.org/licenses/by/4.0/}. Reuse and typesetting adaptation; original language and mathematical content preserved.\par
\noindent\textbf{Difficulty (prior selection preserved):} {\CJKfont H2: The modular representation decomposition interacts with Frobenius powers and changing quasi-invariance thresholds. The proof constructs explicit free generators across the powers-of-two transition and proves completeness through a representation-sensitive induction, rather than checking that the displayed polynomials satisfy the defining divisibility condition.}\par
\medskip\noindent\textbf{Editorial boundary note (not author text):} Complete original Theorem3.6, quasi-invariant/field/symmetric-module and half-integer conventions, characteristic2 indecomposable decomposition, trivial-extension and standard-module degree/freeness core, and full author induction across the dyadic boundary. The current-paper characteristic2 identities and reductions are retained. The author explicitly cites identical prior Wang arguments for the converse/divisibility and residual module-generation steps; those citations/analogies remain transparent, and their prior human-authored arguments are included as a separately attributed arXiv2201.06111v2 component with its original scope visible. Characteristic3 constructions and later Ren–Xu classification are excluded. A narrow mathbbm resolver accepts only literal lowercase k, preserves its field meaning/case, and uses the named installed DejaVu Math TeX Gyre double-struck glyph; this font is an explicit portable dependency and all other arguments are rejected. Unused pgfplots imports are omitted; no plot occurs in the selected source.\par
\medskip\hrule\medskip\noindent\textbf{Author text begins below.}\par
\setcounter{section}{1}
% Original source lines 73--83

Let $\k$ be a field, and consider the action of the symmetric group $S_n$ on the space $\k[x_1,\dots,x_n]$ of $\k$-valued polynomials by permuting the variables. A polynomial in $\k[x_1,\dots,x_n]$ is \emph{symmetric} if it is invariant under this action. Equivalently, since $S_n$ is generated by transpositions, a~polynomial~$K$ is symmetric if $s_{i_1i_2}K=K$ or $(1-s_{i_1i_2})K=0$ for all $s_{i_1i_2}\in S_n$. One may consider generalizations of symmetric polynomials in which this condition is relaxed, so that we only require $(1-s_{i_1i_2})K$ be divisible by some large polynomial. This leads to the notion of \emph{quasi-invariant polynomials}.

\begin{Definition}
 Let $\k$ be a field. For $m\in\Z_{\geq 0}$, $n\in\Z_{>0}$, a polynomial $K\in\k[x_1,\dots,x_n]$ is \textit{$m$-quasi-invariant} if for all $s_{i_1i_2}\in S_n$ we have that $(x_{i_1}-x_{i_2})^{2m+1}$ divides $(1-s_{i_1i_2})K$. We denote the space of $m$-quasi-invariants by $Q_m(n,\k)$.
\end{Definition}

Note that the symmetric polynomials are exactly the polynomials that are $m$-quasi-invariant for all $m$. For brevity, we also refer to quasi-invariant polynomials as simply quasi-invariants.

Quasi-invariant polynomials were first introduced by Chalykh and Veselov in 1990 \cite{cmp/1104179957} to describe the harmonic, zero eigenvalue eigenfunctions of quantum Calogero--Moser systems. Calogero--Moser systems are a collection of one-dimensional dynamical particle systems that were found to be both solvable \cite{calogero} and integrable \cite{MOSER1975197}. Due to these properties, they have become extensively studied in mathematical physics, with connections to a number of other fields of mathematics, including representation theory.


\setcounter{Theorem}{2}
% Original source lines 99--103
\begin{Theorem}\label{thm:1}
 Let $\k$ be either $\mathbf{F}_2$ or $\mathbf{F}_3$. Then the Hilbert series for $Q_m(3,\k)$ is given by
 \[\mathcal{H}(Q_m(3,\k))=\frac{1+2t^d+2t^{6m+3-d}+t^{6m+3}}{(1-t)(1-t^2)(1-t^3)},\]
 where $d=3m+1$ if there is no Ren--Xu counterexample and $d$ is the degree of the minimal degree Ren--Xu counterexample otherwise. In particular, the conditions found in {\rm \cite{Ren_2020}} for the Hilbert series of $Q_m(3,\k)$ to be different from the Hilbert series of $Q_m(3,\mathbf{Q})$ are necessary.
\end{Theorem}


% Original source lines 111--168
\section{Preliminaries}\label{prelim}
We start with some useful properties of the quasi-invariants.
\begin{Proposition}[\cite{etingof2002lectures}]\label{firstProp}
 Let $\k$ be a field.
\begin{enumerate}\itemsep=0pt
 \item[$1$.] $\k[x_1,x_2,x_3]^{S_3}\subset Q_m(3,\k)$, $Q_0(3,\k) = \k[x_1,x_2,x_3]$, and $Q_m(3,\k)\supset Q_{m'}(3,\k)$, where ${m'>m}$.
 \item[$2$.] $Q_m(3,\k)$ is a ring.
 \item[$3$.] $Q_m(3,\k)$ is a finitely generated $\k[x_1,x_2,x_3]^{S_3}$-module.
\end{enumerate}
\end{Proposition}

Note that \cite{etingof2002lectures} proves Proposition~\ref{firstProp} in the case, where $\k=\C$. However, the proofs for the first two assertions work over any field, and the last assertion follows from the Hilbert basis theorem. In view of the structure of $Q_m(3,\k)$ as a module over the symmetric polynomials, given some $K\in Q_m(3,\k)$, we will frequently refer to quasi-invariant polynomials that can be obtained via scalar multiplication of $Q$ by a symmetric polynomial. To distinguish these polynomials from the ordinary $\k$-multiples of $Q$, we will refer to them as \emph{symmetric polynomial multiples} of $Q$.

We consider $Q_m(3, \mathbf{F}_{2})$ and $Q_m(3, \mathbf{F}_{3})$ as representations of $S_3$, where $S_3$ permutes the variables $x_1$, $x_2$, $x_3$. Since $Q_m(3, \mathbf{F}_{2})$ and $Q_m(3, \mathbf{F}_{3})$ are vector spaces over $\mathbf{F}_2$ and $\mathbf{F}_3$ respectively and the characteristics $2$ and $3$ divide $|S_3|$, $Q_m(3, \mathbf{F}_{2})$ and $Q_m(3, \mathbf{F}_{3})$ are modular representations~of~$S_3$.

\begin{Proposition}
 $Q_m(3, \mathbf{F}_{2})$ and $Q_m(3, \mathbf{F}_{3})$ are modular representations of $S_3$.
\end{Proposition}

First, we consider characteristic $2$.

\subsection[Preliminary definitions for p=2]{Preliminary definitions for $\boldsymbol{p=2}$}

We describe the indecomposable and irreducible representations of $S_3$ for $p=2$.

\begin{Proposition}[\cite{Alperin_1986}]\label{prop:p2indecomp}
 There are $3$ irreducible or indecomposable representations of $S_{3}$ in characteristic $2$:
 \begin{enumerate}\itemsep=0pt
 \item[$1$.] $\triv$ is the irreducible representation of $S_3$ that is acted on trivially by $S_{3}$.
 \item[$2$.] $\std$ is the $2$-dimensional irreducible representation of $S_3$ obtained by reducing the standard representation in characteristic $0$ mod~$2$.
 \item[$3$.] $\trtr$ is the $2$-dimensional indecomposable representation that contains a copy of $\triv$ as a subrepresentation such that the quotient of $\trtr$ by this subrepresentation is $\triv$.
 \end{enumerate}
\end{Proposition}

\begin{Example}\label{eg:p2reps}
 The polynomial $E_{\trtr}:=x_1^2x_2+x_2^2x_3+x_3^2x_1\in\mathbf{F}_2[x_1,x_2,x_3]$ generates a~copy of $\trtr$. To see this, note that for any $i_1$, $i_2$, we have
 \[(1-s_{i_1i_2})E_\trtr=x_1^2x_2+x_1x_2^2+x_1^2x_3+x_1x_3^2+x_2^2x_3+x_2x_3^2\in\mathbf{F}_2[x_1,x_2,x_3]^{S_3}.\]
 Since the transpositions generate $S_3$, $E_\trtr$ generates a two-dimensional representation that contains $\triv$ as a subrepresentation. Moreover, since $E_\trtr$ is not symmetric, this representation is not $\triv\oplus\triv$, so it must be $\trtr$.
\end{Example}

We then study the behaviors of each indecomposable representation in the quasi-invariants. We define $Q_m(3,\mathbf{F}_2)_{\triv}$ and $Q_m(3,\mathbf{F}_2)_{\std}$ to be the direct sum of all copies of $\triv$ and $\std$ respectively in the quasi-invariants. We also define $Q_m(3,\mathbf{F}_2)_\trtr$ to be the direct sum of all copies of $\triv$ and $\trtr$.

\begin{Remark}
 We cannot define $Q_m(3,\mathbf{F}_2)_\trtr$ to exclude copies of $\triv$ since we can add elements of $Q_m(3,\mathbf{F}_2)_\triv$ to copies of $\trtr$ and still obtain a copy of $\trtr$. For example, \smash{$F:=E_\trtr+x_1^3+x_2^3+x_3^3$} still satisfies $(1-s_{i_1i_2})F=(1-s_{i_1i_2})E_\trtr$ for all~$i_1$,~$i_2$, so it generates a copy of $\trtr$ by the same argument as Example~\ref{eg:p2reps}.
\end{Remark}

\begin{Proposition}[{\cite{wang2022explicit}}]\label{prop:p2triv}
 As an $\mathbf{F}_2[x_1,x_2,x_3]^{S_3}$-module, $Q_m(3,\mathbf{F}_2)_{\triv}$ is freely generated by $1$.
\end{Proposition}

Note that by the classification of indecomposables in Proposition~\ref{prop:p2indecomp}, every extension of $\std$ and every extension of a module by $\std$ splits. Thus $Q_m(3,\mathbf{F}_2)_\std$ is a direct summand of $Q_m(3,\mathbf{F}_2)$ (whose complement is $Q_m(3,\mathbf{F}_2)_\trtr$), and we mainly consider $Q_m(3,\mathbf{F}_2)_\std$. $Q_m(3,\mathbf{F}_2)_{\std}$ is generated as a $\mathbf{F}_2[x_1,x_2,x_3]^{S_3}$-module by homogeneous copies of $\std$, so following~\cite{wang2022explicit}, we consider \textit{generating representations} of $Q_m(3,\mathbf{F}_2)_\std$ as homogeneous copies of $\std$ in a generators and relations presentation of $Q_m(3,\mathbf{F}_2)_{\std}$ with a minimal generator set.

\subsubsection{Quasi-invariants at half-integers}

Note that if $\k$ is a field with $\mathrm{char\,}\k\neq 2$ and $m\in\Z_{\geq 0}$, then for any $K\in\k[x_1,\dots,x_n]$, $(x_{i_1}-x_{i_2})^{2m}|(1-s_{i_1i_2})K$ implies $(x_{i_1}-x_{i_2})^{2m+1}|(1-s_{i_1i_2})K$ since $(1-s_{i_1i_2})K$ is $s_{i_1i_2}$-antiinvariant, hence the exponent $2m+1$ in the definition of quasi-invariant polynomials. But this does not hold in characteristic 2, since there is no concept of antiinvariants. Indeed, one can check that for $K=x_1^2+x_2^2$, we have $(x_{i_1}-x_{i_2})^2|(1-s_{i_1i_2})K$ for all $i_1$, $i_2$, but $(x_{i_1}-x_{i_2})^3\nmid|(1-s_{i_1i_2})K$ if~$i_1=1,2$,~$i_2\neq 1,2$.

We encapsulate this data by extending the definition of quasi-invariants to half-integers when ${p=2}$. For example, $K=x_1^2+x_2^2$ is \smash{$\frac{1}{2}$}-quasi-invariant, and this is in fact the minimal degree nonsymmetric \smash{$\frac{1}{2}$}-quasi-invariant polynomial. Proposition~\ref{firstProp} still holds when $m,m'$ are half-integers, and the definitions of $Q_m(3,\mathbf{F}_2)_\triv$, $Q_m(3,\mathbf{F}_2)_\std$ also naturally extend to half-integer~$m$. So from now on, whenever we refer to quasi-invariants in characteristic $2$ we let $m$ be a~half-integer.


\setcounter{section}{2}
% Original source lines 227--306
\section[Quasi-invariants in characteristic 2]{Quasi-invariants in characteristic 2}\label{sec:p=2}

In this section, we write down explicit generators for $Q_m(3,\mathbf{F}_2)$ and prove Theorem~\ref{thm:1} for $p=2$. Note that we already know the structure of $Q_m(3,\mathbf{F}_2)_\triv$ from Proposition~\ref{prop:p2triv}. We start by extending this to $Q_m(3,\mathbf{F}_2)_\trtr$.

\begin{Proposition}\label{prop:trivtriv}
 As an $\mathbf{F}_2[x_1,x_2,x_3]^{S_3}$-module, $Q_m(3,\mathbf{F}_2)_{\triv-\triv}$ is freely generated by $1$ and $E_{\trtr}\prod(x_{i_1}-x_{i_2})^{2m}$.
\end{Proposition}
\begin{proof}
 Let $K$ be a nonsymmetric element of $Q_m(3,\mathbf{F}_2)_{\triv-\triv}$ so that $(x_{i_1}-x_{i_2})^{2m+1}$ divides $(1+s_{i_1i_2})K$. Because
 \[(1+s_{12})K=(1+s_{13})K=(1+s_{23})K,\]
 we have $(1+s_{i_1i_2})K = P\prod (x_{i_1}-x_{i_2})^{2m+1}$ for some symmetric polynomial $P$. Letting $G=E_{\trtr}\prod (x_{i_1}-x_{i_2})^{2m}$ yields $(1+s_{i_1i_2})G = \prod (x_{i_1}-x_{i_2})^{2m+1}$. Thus $(1+s_{i_1i_2})PG = (1+s_{i_1i_2})K$ and $(1+s_{i_1i_2})(PG-K)=0$, so $PG-K$ is symmetric and $K$ is generated by $G$ and $1$. Moreover, since $G$ is not symmetric, $P$ and $G$ have no relation implying freeness.
\end{proof}

We have an explicit description of $Q_m(3,\mathbf{F}_2)_\trtr$, so it remains to compute the generators and relations of $Q_m(3,\mathbf{F}_2)_\std$. A number of the properties of $Q_m(3,\mathbf{F}_p)$ for $p>3$ found in \cite{wang2022explicit} are true for $Q_m(3,\mathbf{F}_2)$. We prove these first.

If $V$ is a copy of $\std$, then we denote by $V_{i_1i_2}$ the $1$-eigenspace of $s_{i_1i_2}$ in $V$.

\begin{Lemma}
 Let $V$ be a copy of $\std$ in $Q_m(3,\mathbf{F}_2)_{\std}$, and let $K\in V_{i_1i_2}$. Then we have ${K+sK+s^2K}=0$, where $s=(1\,2\,3)\in S_3$ and $K=(x_{i_1}-x_{i_2})^{2m+1}K'$ for some polynomial~$K'$ that is invariant under the action of $s_{i_1i_2}$. Conversely, let $K'$ be an $s_{12}$-invariant polynomial such that
 \[(x_1-x_2)^{2m+1}K'+(x_2-x_3)^{2m+1}sK'+(x_3-x_1)^{2m+1}s^2K'=0.\]
 Then $(x_1-x_2)^{2m+1}K'$ belongs to the $1$-eigenspace of $s_{12}$ in some copy of $\std$ inside $Q_m(3,\mathbf{F}_2)_\std$.
\end{Lemma}
\begin{proof}
 For the first statement, $K+sK+s^2K=0$ holds for any copy of $\std$. For the next, suppose $\{i_1,i_2,i_3\} = \{1,2,3\}$ for some integer $i_3$. Then $(1-s_{i_1i_3})K = s_{i_2i_3}K$, so $(x_{i_1}-x_{i_3})^{2m+1}|s_{i_2i_3}K$, implying $(x_{i_1}-x_{i_2})^{2m+1}|K$.
 The second statement follows from the proof in~\cite{wang2022explicit}.
\end{proof}

\begin{Corollary}
 Let $V$ be a generating representation of $Q_m(3,\k)_\std$ and let $K\in V_{i_1i_2}$. Let us write $K=(x_{i_1}-x_{i_2})^{2m+1}K'$. Then $K'$ is not divisible by any nonconstant symmetric polynomial.
\end{Corollary}

The proof of this statement is identical to the one in \cite{wang2022explicit}.

\begin{Lemma}\label{char3.4}
 Let $V$, $W$ be distinct generating representations of $Q_m(3,\k)_\std$. Let ${K\in V_{12}}$, ${L\in W_{12}}$. For $\sigma K\sigma L:=(\sigma K)(\sigma L)$, we have that $KL+s_{13}Ks_{23}L$ is a nonsymmetric element of $Q_m(3,\k)_\trtr$ and $\deg V+\deg W\geq 6m+3$.
\end{Lemma}
\begin{proof}
 $KL+s_{13}Ks_{23}L$ is an element of $Q_m(3,\mathbf{F}_2)$ since the quasi-invariants form a ring by Proposition~\ref{firstProp}. Using that $s_{12}K=K$ and $s_{12}L=L$, we have that
 \begin{gather*}
 (1+s_{12})(KL+s_{13}Ks_{23}L) = s_{23}Ks_{13}L+s_{13}Ks_{23}L,\\
 (1+s_{13})(KL+s_{13}Ks_{23}L) = KL+s_{13}Ks_{23}L+s_{13}Ks_{13}L + Ks_{23}L = Ks_{13}L+s_{13}KL,\\
 (1+s_{23})(KL+s_{13}Ks_{23}L) = KL+s_{13}Ks_{23}L+s_{23}Ks_{23}L + s_{13}KL = Ks_{23}L+s_{23}KL.
 \end{gather*}
 One can check that each polynomial is a transposition of another and that they are symmetric due to the structure of $\trtr$, so they are all the same symmetric polynomial. Thus ${KL+s_{13}Ks_{23}L}$ lies in a quotient of a copy of $\trtr$. Note that by the same argument as in \cite{wang2022explicit}, we have $Ks_{23}L+s_{23}KL\neq 0$, so $KL+s_{13}Ks_{23}L$ is nonsymmetric and must generate a~copy of $\trtr$.

 By Proposition~\ref{prop:trivtriv}, $KL+s_{13}Ks_{23}L$ has degree at least $6m+3$, so $\deg V+\deg W\geq 6m+3$ as desired.
\end{proof}

\begin{Lemma}\label{charFree}
 Assume that there exist generating representations $V$, $W$ of $Q_m(3,\mathbf{F}_2)_\std$ such that $\deg V+\deg W=6m+3$. Then $Q_m(3,\mathbf{F}_2)_\std$ is a free module over $\k[x_1,x_2,x_3]^{S_3}$ generated by~$V$ and $W$.
\end{Lemma}
\begin{proof}
 Assume for the sake of contradiction there exists another generator $U$ of $Q_m(3,\mathbf{F}_2)_{\std}$. Supposing $\deg W\geq \deg V$, by Lemma~\ref{char3.4}, $\deg U\geq \deg W$. By Lemma~\ref{char3.4}, if $K\in V_{12}$, $L\in W_{12}$, and $T\in U_{12}$ then $KL+s_{13}Ks_{23}L$ and $KT+s_{13}Ks_{23}T$ are both nonsymmetric elements of $Q_m(3,\mathbf{F}_2)_{\trtr}$. Moreover, we have
 \[(1+s_{12})(KL+s_{13}Ks_{23}L) = s_{23}Ks_{13}L+s_{13}Ks_{23}L = \prod (x_{i_1}-x_{i_2})^{2m+1},\]
 and
 \[(1+s_{12})(KT+s_{13}Ks_{23}T) = s_{23}Ks_{13}T+s_{13}Ks_{23}T = Q\prod (x_{i_1}-x_{i_2})^{2m+1}\]
 for some symmetric polynomial $Q$. From there we may proceed identically to \cite{wang2022explicit}.
\end{proof}

Now, we are ready to prove Theorem~\ref{thm:1} for $p=2$.

\begin{Theorem}\label{thm:p2}
 Let $a$ be the largest natural number such that $2^a< 2m+1$. Then $Q_m(3,\mathbf{F}_2)_{\std}$ is freely generated by $(x_1-x_2)^{2^{a+1}}$ and $(x_1-x_2)^{2^a}\prod (x_{i_1}-x_{i_2})^{2m+1-2^a}$.
\end{Theorem}

\begin{Remark}
 Note that when $m$ is an integer, the degrees of the generators in this theorem agree with the degrees conjectured in \cite{Ren_2020}. In particular, when $2^{a+1}$ is one of $3m+1$, ${3m+2}$, we actually have that the Hilbert series of $Q_m(3,\F_2)$ and $Q_m(3,\mathbf{Q})$ agree, so $(x_1-x_2)^{2^{a+1}}$, \smash{$(x_1-x_2)^{2^a}\prod (x_{i_1}-x_{i_2})^{2m+1-2^a}$} are the reductions modulo 2 of the generators of $Q_m(3,\mathbf{Q})$, when written as integer polynomials with coprime coefficients.
\end{Remark}

\begin{proof}[Proof of Theorem~\ref{thm:p2}]
 We prove this by induction on $m$.

 The generators of $Q_0(3,\mathbf{F}_2)_{\std}$ are $(x_1-x_2)$ and $(x_1-x_2)^2$, completing our base case.

 Let $j$ be a half-integer, and suppose that \smash{$Q_{j-\frac{1}{2}}(3,\mathbf{F}_2)_\std$} is freely generated by $(x_1-x_2)^{2^{a+1}}$ and \smash{$(x_1-x_2)^{2^a}\prod (x_{i_1}-x_{i_2})^{2j-2^a}$},
 where $2^a$ is the greatest such power of $2$ less than $2j$. If~${2j\neq 2^{a+1}}$, then $2^{a}$ is the largest power of $2$ less than $2j+1$, so $(x_1-x_2)^{2^{a+1}}$ and $(x_1-x_2)^{2^a}\prod (x_{i_1}-x_{i_2})^{2j+1-2^a}$ are both in $Q_j(3,\mathbf{F}_2)$. Further, $(x_1-x_2)^{2^{a+1}}$ must be a generator and if $(x_1-x_2)^{2^a}\prod (x_{i_1}-x_{i_2})^{2j+1-2^a}$ is a not a generator, by Lemma~\ref{char3.4}, $(x_1-x_2)^{2^a}\prod (x_{i_1}-x_{i_2})^{2j+1-2^a}$ is generated by $(x_1-x_2)^{2^{a+1}}$ which implies a relation between $(x_1-x_2)^{2^{a+1}}$ and ${(x_1-x_2)^{2^a}\prod (x_{i_1}-x_{i_2})^{2j-2^a}}$. Because they freely generate \smash{$Q_{j-\frac{1}{2}}(3,\mathbf{F}_2)$}, this is impossible. Thus $(x_1-x_2)^{2^{a+1}}$ and $(x_1-x_2)^{2^a}\smash{\prod (x_{i_1}-x_{i_2})^{2j+1-2^a}}$ freely generate $Q_j(3,\mathbf{F}_2)_{\std}$ by Lemma~\ref{charFree}.

 If $2j= 2^{a+1}$, then both $(x_2-x_3)^{2^{a+1}}$ and $(x_2-x_3)^{2^{a+2}}\prod (x_{i_1}-x_{i_2})^{2j+1-2^a}$ lie in $Q_j(3,\mathbf{F}_2)_{\std}$. The former is a generator by our inductive hypothesis.
 Since $2^{a+1}+2^{a+2}+3=6j+3$, if the latter is not a generator, then by Lemma~\ref{char3.4}, $(x_2-x_3)^{2^{a+2}\prod (x_{i_1}-x_{i_2})^{2j+1-2^a}}$ is generated by $(x_2-x_3)^{2^{a+1}}$, which is false. Thus $(x_2-x_3)^{2^{a+1}}$ and $(x_2-x_3)^{2^{a+2}}\prod (x_{i_1}-x_{i_2})^{2j+1-2^a}$ freely generate $Q_j(3,\mathbf{F}_2)_{\std}$ by Lemma~\ref{charFree} as desired.
\end{proof}
\clearpage
\section*{Separately attributed prior human-authored proof support}
\noindent\textbf{Author:} Frank Wang. \textbf{Work:} Toward explicit Hilbert series of quasi-invariant polynomials in characteristic $p$ and $q$-deformed quasi-invariants.\par
\noindent\textbf{Exact component version:} arXiv:2201.06111v2, submitted 29 September 2022, \url{https://arxiv.org/abs/2201.06111v2}. \textbf{Rights:} CC BY4.0, \url{https://creativecommons.org/licenses/by/4.0/}. Author text retained; standalone formatting, cross-reference namespace and bibliography presentation adapted.\par
\noindent\textbf{Editorial source boundary:} This is the preprint-v2 proof source, separately identified from the NYJM29(2023),613--634 bibliographic identity cited by the 2025 paper. No version identity is asserted. The prior non-modular and characteristic $p>3$ hypotheses remain visible; no hypothesis is transferred to characteristic2 or3 by the editor. The 2025 authors' displayed characteristic-specific identities and exact analogies remain in their own text above. This component supplies the prior Lemma3.2 converse, Corollary3.3, Lemma3.4 and Lemma3.5 author arguments with their original setting. The mention of prior Section4 is a preview, not an omitted part of these proofs.\par
\medskip\hrule\medskip
\begingroup
\theoremstyle{plain}
\newtheorem{theorem}{Theorem}[section]
\newtheorem{corollary}[theorem]{Corollary}
\newtheorem{lemma}[theorem]{Lemma}
\newtheorem{proposition}[theorem]{Proposition}
\newtheorem{conjecture}[theorem]{Conjecture}

\theoremstyle{definition}
\newtheorem{definition}[theorem]{Definition}
\newtheorem{example}[theorem]{Example}
\newtheorem{remark}[theorem]{Remark}

\setcounter{section}{1}
\let\WangOriginalLabel\label
\let\WangOriginalRef\ref
\let\WangOriginalEqref\eqref
\let\WangOriginalCite\cite
\makeatletter\let\WangOriginalDisplayLabel\label@in@display\def\label@in@display#1{\WangOriginalDisplayLabel{wangv2:#1}}\makeatother
\renewcommand{\label}[1]{\WangOriginalLabel{wangv2:#1}}
\renewcommand{\ref}[1]{\WangOriginalRef{wangv2:#1}}
\renewcommand{\eqref}[1]{\WangOriginalEqref{#1}}
\renewcommand{\cite}[1]{\WangOriginalCite{wangv2:#1}}
\renewcommand{\char}{\text{char\,}}
\setlist[enumerate]{label=\arabic*),ref=\arabic*)}
% Original Wang arXiv v2 source lines81--249
\section{Quasi-invariant polynomials for $S_n$}\label{sec:2}

We start with the definition of quasi-invariant polynomials. Let $\k$ be a field, and let $S_n$ act on $\k[x_1,\dots,x_n]$ by permuting the variables. Let $s_{i,j}\in S_n$ denote the transposition swapping the $i$-th and $j$-th indices.

\begin{definition}
    Let $n$ be a positive integer. A polynomial $P\in\k[x_1,\dots,x_n]$ is $m$-quasi-invariant if for all transpositions $s_{i,j}\in S_n$, $(1-s_{i,j})P$ is divisible by $(x_i-x_j)^{2m+1}$. We denote by $Q_m(n,\k)$ the space of all $m$-quasi-invariants in $\k[x_1,\dots,x_n]$. 
\end{definition}

Let $\k[x_1,\dots,x_n]^{S_n}$ denote the space of polynomials in $\k[x_1,\dots,x_n]$ which are invariant with respect to the action of $S_n$. Some useful properties of $Q_m(n,\k)$ are described in the following proposition. 

\begin{proposition}[\cite{etingof2002lectures}]\label{prop:2.2}
    \phantom{.}

    \begin{enumerate}
        \item $\k[x_1,\dots,x_n]^{S_n}\subset Q_m(n,\k)$, $Q_0(n,\k)=\k[x_1,\dots,x_n]$, $Q_m(n,\k)\subset Q_{m'}(n,\k)$ if $m>m'$.
        \item $Q_m(n,\k)$ is a ring.
        \item $Q_m(n,\k)$ is a finitely generated module over $\k[x_1,\dots,x_n]^{S_n}$.
    \end{enumerate}
\end{proposition}

Note that $Q_m(n,\k)$ admits a grading by the degrees of the polynomials. Thus we may define its Hilbert series.

\begin{definition}
    The Hilbert series of $Q_m(n,\k)$ is
    $$H_m(t)=\sum_{d\geq 0}\dim Q_m(n,\k)[d]t^d$$
    where $Q_m(n,\k)[d]$ is the graded component of $Q_m(n,\k)$ consisting of polynomials of homogeneous degree $d$.
\end{definition}

Note that by the fundamental theorem of symmetric polynomials, $\k[x_1,\dots,x_n]^{S_n}$ is a free polynomial algebra generated by symmetric polynomials $P_1,\dots,P_n$ of degrees $1,\dots,n$, respectively \cite{humphreys_1990}. Thus since $Q_m(n,\k)$ is a finitely generated module over $\k[x_1,\dots,x_n]^{S_n}=\k[P_1,\dots,P_n]$ by Proposition \ref{prop:2.2}, we can write
$$H_m(t)=\frac{G_m(t)}{\prod_{d=1}^n(1-t^d)}$$
where $G_m(t)$ is a polynomial by Hilbert's syzygy theorem. We say that $G_m(t)$ is the Hilbert polynomial associated with $H_m(t)$. Note that describing the Hilbert polynomial is related to describing the generators and relations of $Q_m(n,\k)$ as a $\k[x_1,\dots,x_n]^{S_n}$-module.

From now on, we will assume the characteristic of $\k$ does not divide $n!$, so that the representation theory of $S_n$ over $\k$ is non-modular. The following lemma and proposition are standard results in the study of quasi-invariants, for example they a direct consequence of the decomposition of $Q_m(n,\C)$ as a representation of the spherical rational Cherednik algebra (see \cite{berest2003cherednik}) in the case $\k=\C$. 

\begin{lemma}\label{lem:2.4}
    $Q_m(n,\k)$ is a representation of $S_n$.
\end{lemma}

Now, consider a graded component $Q_m(n,\k)[d]$ of the quasi-invariants. It is a subrepresentation of $Q_m(n,\k)$, and moreover it is finite dimensional since it is a subspace of the finite dimensional space $\k[x_1,\dots,x_n][d]$. Thus, since we are working in the non-modular case, by Maschke's theorem it decomposes as a direct sum
$$Q_m(n,\k)[d]=\bigoplus_{\tau\in\Irrep(S_n)}Q_m(n,\k)[d]_\tau$$
where $\Irrep(S_n)$ is the set of irreducible representations of $S_n$ and $Q_m(n,\k)[d]_\tau$ is the subspace of $Q_m(n,\k)[d]$ on which $S_n$ acts by a direct sum of copies of $\tau$. Then for $\tau\in\Irrep(S_n)$, let us define
$$Q_m(n,\k)_\tau=\bigoplus_{d\geq 0}Q_m(n,\k)[d]_\tau$$
to be the isotypic components of $Q_m(n,\k)$.

\begin{proposition}\label{prop:2.5}
\begin{enumerate}
    \item $Q_m(n,\k)_\tau$ is a module over $\k[x_1,\dots,x_n]^{S_n}$.\label{prop:2.51}
    \item We have the decomposition
    $$Q_m(n,\k)=\bigoplus_{\tau\in\Irrep(S_n)}Q_m(n,\k)_\tau$$
    as a $\k$-vector space and hence as a $\k[x_1,\dots,x_n]^{S_n}$-module as well.\label{prop:2.52}
\end{enumerate}
\end{proposition}

By the above proposition, to study the generators and relations of $Q_m(n,\k)$ it suffices to study the generators and relations of $Q_m(n,\k)_\tau$. This will be our primary method of studying $Q_m(n,\k)$ throughout the rest of this paper. Let us denote by triv and sign the trivial and sign representations of $S_n$, respectively. We deal with these two cases here.

\begin{proposition}\label{prop:2.6} Assume $\char\k\neq 2$. As $\k[x_1,\dots,x_n]^{S_n}$-modules, we have
\begin{enumerate}
    \item $Q_m(n,\k)_\triv$ is freely generated by $1$.\label{prop:2.61}
    \item $Q_m(n,\k)_\sgn$ is freely generated by $\prod_{i<j}(x_i-x_j)^{2m+1}$.\label{prop:2.62}
\end{enumerate}
\end{proposition}
\begin{proof}
    \ref{prop:2.61}~Note that the statement that $S_n$ acts by the trivial representation on $P$ is equivalent to the statement that $P$ is $S_n$-invariant, so $Q_m(n,\k)_\triv\subset\k[x_1,\dots,x_n]^{S_n}$. On the other hand, for any $P\in\k[x_1,\dots,x_n]^{S_n}$ and any $s_{i,j}\in S_n$ we have $(1-s_{i,j})P=P-P=0$ is divisible by $(x_i-x_j)^{2m+1}$, so $Q_m(n,\k)_\triv=\k[x_1,\dots,x_n]^{S_n}$.
    
    \ref{prop:2.62}~Let $P\in Q_m(n,\k)_\sgn$ and let $s_{i,j}\in S_n$. Then since $P$ is in the sign representation we have $s_{i,j}P=-P$, so $(1-s_{i,j})P=2P$ is divisible by $(x_i-x_j)^{2m+1}$. So $P$ is divisible by $(x_i-x_j)^{2m+1}$ for all $i$ and $j$, and thus it is divisible by $\prod_{i<j}(x_i-x_j)^{2m+1}$. 
    
    Let us write $P=K\prod_{i<j}(x_i-x_j)^{2m+1}$ for some polynomial $K$. Note that 
    $$s_{i',j'}\prod_{i<j}(x_i-x_j)^{2m+1}=-\prod_{i<j}(x_i-x_j)^{2m+1}$$
    for all $s_{i',j'}\in S_n$, so we have 
    $$s_{i',j'}P=-P=-K\prod_{i<j}(x_i-x_j)^{2m+1}=s_{i',j'}K\bigg(s_{i',j'}\prod_{i<j}(x_i-x_j)^{2m+1}\bigg)=-s_{i',j'}K\prod_{i<j}(x_i-x_j)^{2m+1}$$
    so $s_{i',j'}K=K$. Since $S_n$ is generated by transpositions, $K$ is $S_n$-invariant, and thus $P$ is in the $\k[x_1,\dots,x_n]^{S_n}$-module generated by $\prod_{i<j}(x_i-x_j)^{2m+1}$. Conversely, for any $K\in\k[x_1,\dots,x_n]^{S_n}$ and $s_{i',j'}\in S_n$, we have 
    $$s_{i',j'}\Big(K\prod_{i<j}(x_i-x_j)^{2m+1}\Big)=-K\prod_{i<j}(x_i-x_j)^{2m+1}$$
    and $$(1-s_{i',j'})\Big(K\prod_{i<j}(x_i-x_j)^{2m+1}\Big)=2K\prod_{i<j}(x_i-x_j)^{2m+1}$$ is divisible by $(x_{i'}-x_{j'})^{2m+1}$, so the $\k[x_1,\dots,x_n]^{S_n}$-module generated by $\prod_{i<j}(x_i-x_j)^{2m+1}$ is contained in $Q_m(n,\k)_\sgn$. So $Q_m(n,\k)_\sgn$ is exactly this module. Freeness then follows from the fact that $\k[x_1,\dots,x_n]$ is an integral domain.
\end{proof}

\begin{remark}\label{rem:2.7}
    It was proved in \cite{Ren_2020} that the Hilbert series of $Q_m(2,\k)$ is $\frac{1+t^{2m+1}}{(1-t)(1-t^2)}$, regardless of the characteristic of $\k$. Proposition \ref{prop:2.6} provides an alternate proof of this result in the case $\char \k\neq 2$. Namely, note that the only irreducible representations of $S_2$ are triv and sign. Thus the generators of $Q_m(2,\k)$ are 1 (of degree 0) and $(x_1-x_2)^{2m+1}$ (of degree $2m+1$). Since these generators belong to different irreducible representations of $S_2$, there are no relations between them, and thus the Hilbert polynomial is $1+t^{2m+1}$ and the Hilbert series is $\frac{1+t^{2m+1}}{(1-t)(1-t^2)}$.
\end{remark}

\section{The Hilbert series of $Q_m(3,\k)$}\label{sec:3}

The Hilbert series of $Q_m(3,\C)$ was computed in \cite{felder2003action} and is as follows:
\begin{equation}\label{eq:3.1}
    H_m(t)=\frac{1+2t^{3m+1}+2t^{3m+2}+t^{6m+3}}{(1-t)(1-t^2)(1-t^3)}.\tag{3.1}
\end{equation}

In this section, we work toward an explicit description of the Hilbert series of $Q_m(3,\k)$ in characteristic $p>3$. Note that there are three irreducible representations of $S_3$: triv, sign, and the 2-dimensional standard representation which we will denote by std. Proposition \ref{prop:2.6} computes the generators of $Q_m(3,\k)_\triv$ and $Q_m(3,\k)_\sgn$, so to compute the Hilbert series of $Q_m(3,\k)$ we only need to consider the space $Q_m(3,\k)_\std$. Both this section and the next section will be largely focused on understanding this space.

This section is dedicated to proving the main result of this paper, which is the following theorem, conjectured in \cite{Ren_2020}. Note the similarity to (\ref{eq:3.1}).

\begin{theorem}\label{thm:3.1}
    Let $\k$ be a field of characteristic $p>3$. Then $Q_m(3,\k)_\std$ is freely generated by standard representations of degree $d$ and $6m+3-d$ where $d$ is an integer depending on $m,p$ that satisfies $2m+1\leq d\leq 3m+1$. In particular, the Hilbert series of $Q_m(3,\k)$ is
    $$H_m(t)=\frac{1+2t^d+2t^{6m+3-d}+t^{6m+3}}{(1-t)(1-t^2)(1-t^3)}.$$
    
\end{theorem}

To start, let $V$ be a copy of the standard representation, and let $s_{i,j}\in S_3$. Then $V$ contains a one dimensional $-1$ eigenspace of $s_{i,j}$. Let us denote this eigenspace by $V_{i,j}^-$. Then we have

\begin{lemma}\label{lem:3.2}
    Let $V\subset Q_m(3,\k)_\std$ be a copy of the standard representation, and let $P\in V_{i,j}^-$. Then we have $P+sP+s^2P$=0 where $s=(1\,2\,3)\in S_3$ and $P=(x_i-x_j)^{2m+1}K$ for some polynomial $K$ that is invariant under the action of $s_{i,j}$. Conversely, let $K$ be an $s_{1,2}$-invariant polynomial such that 
    $$(x_1-x_2)^{2m+1}K+(x_2-x_3)^{2m+1}sK+(x_3-x_1)^{2m+1}s^2K=0.$$
    Then $(x_1-x_2)^{2m+1}K$ either belongs to $Q_m(3,\k)_\sgn$ or the $-1$ eigenspace of $s_{1,2}$ in some copy of std inside $Q_m(3,\k)_\std$.
\end{lemma}
\begin{proof}
    The proof of the first statement is analogous to the proof of Proposition \ref{prop:2.6}. Since $P\in V_{i,j}^-$, we have $(1-s_{i,j})P=2P$ is divisible by $(x_i-x_j)^{2m+1}$, so $P=(x_i-x_j)^{2m+1}K$ for some polynomial $K$. Then since $P$ and $(x_i-x_j)^{2m+1}$ are both $s_{i,j}$-antiinvariant, $K$ must be $s_{i,j}$-invariant. The fact that $P+sP+s^2P=0$ is a direct consequence of $P$ being an element of the standard representation.
    
    For the second statement, note that since $K$ is $s_{1,2}$-invariant, $sK=s_{1,3}K$ and $s^2K=s_{2,3}K$. So we have
    $$(1-s_{1,3})\Big((x_1-x_2)^{2m+1}K\Big)=(x_1-x_2)^{2m+1}K-(x_3-x_2)^{2m+1}sK=-(x_3-x_1)^{2m+1}s^2K,$$
    $$(1-s_{2,3})\Big((x_1-x_2)^{2m+1}K\Big)=(x_1-x_2)^{2m+1}K-(x_1-x_3)^{2m+1}s^2K=-(x_2-x_3)^{2m+1}sK$$
    so $(x_1-x_2)^{2m+1}K$ is $m$-quasi-invariant. On the other hand, since $s_{1,2}$ sends the vector space spanned by $(x_1-x_2)^{2m+1}K$ to itself and $s_{1,3}$ and $s_{2,3}$ send $(x_1-x_2)^{2m+1}K$ to polynomials in the space spanned by $(x_1-x_2)^{2m+1}K$ and $(x_2-x_3)^{2m+1}sK$, the span of the $S_3$ orbit of $(x_1-x_2)^{2m+1}K$ is at most 2 dimensional. If it is 1 dimensional, then $(x_1-x_2)^{2m+1}K$ is contained in a copy of sign inside $Q_m(3,\k)_\sgn$ since $s_{1,2}$ acts by $-1$. If it is 2 dimensional, then since the only 2 dimensional irreducible representation of $S_3$ is std, $(x_1-x_2)^{2m+1}K$ belongs to some copy of std inside $Q_m(3,\k)_\std$. That it is in the $-1$ eigenspace of $s_{1,2}$ is clear by applying $s_{1,2}$ to $(x_1-x_2)^{2m+1}K$.
\end{proof}

Note that since $Q_m(3,\k)_\std$ is a graded $\k[x_1,x_2,x_3]^{S_3}$-module, we may assume its generators are homogeneous. We may also describe generating representations of $Q_m(3,\k)_\std$, rather than generating elements. 

\begin{corollary}\label{cor:3.3}
    Let $V$ be a generating representation of $Q_m(3,\k)_\std$ and let $P\in V_{i,j}^-$. Let us write $P=(x_i-x_j)^{2m+1}K$. Then $K$ is not divisible by any nonconstant symmetric polynomial.
\end{corollary}
\begin{proof}
    Since $P$ is in $Q_m(3,\k)_\std$, it satisfies $P+sP+s^2P=0$ where $s=(1\,2\,3)$. Assume for contradiction that $K$ is divisible by some symmetric polynomial $Q$. Then it is clear that \linebreak
    $P/Q+s(P/Q)+s^2(P/Q)=0$ and $P/Q=(x_i-x_j)^{2m+1}(K/Q)$, so $P/Q$ is contained in $Q_m(3,\k)$ by Lemma \ref{lem:3.2}. But then $P$ is in the $\k[x_1,x_2,x_3]^{S_3}$-module generated by $P/Q$, so it is generated by an element of lower degree and thus is not contained in a generating representation.
\end{proof}

For a standard representation $V$ consisting of homogeneous polynomials of degree $d$, let us say its degree is $\deg V=d$. The following lemma makes use of this, and will prove to be useful both in the proof of Theorem \ref{thm:3.1} and in Section \ref{sec:4}.

\begin{lemma}\label{lem:3.4}
    Let $V,W$ be distinct generating representations of $Q_m(3,\k)_\std$. Let $A\in V_{1,2}^-,B\in W_{1,2}^-$. Then $As_{2,3}B-Bs_{2,3}A$ is a nonzero element of $Q_m(3,\k)_\sgn$ and we have $\deg V+\deg W\geq 6m+3$.
\end{lemma}
\begin{proof}
    $As_{2,3}B-Bs_{2,3}A$ is contained in $Q_m(3,\k)$ because the quasi-invariants form a ring by Proposition \ref{prop:2.2}. It is also contained in a sign representation since it spans the representation $\wedge^2\std=\sgn$ inside $V\otimes W$. Thus it is contained in $Q_m(3,\k)_\sgn$, so it remains to show that it is nonzero.
    
    By Lemma \ref{lem:3.2} we can write $A=(x_1-x_2)^{2m+1}K,B=(x_1-x_2)^{2m+1}L$ for $s_{1,2}$-invariant polynomials $K$ and $L$. Then we have
    \begin{align*}
        As_{2,3}B-Bs_{2,3}A&=(x_1-x_2)^{2m+1}K(x_1-x_3)^{2m+1}s_{2,3}L-(x_1-x_2)^{2m+1}L(x_1-x_3)^{2m+1}s_{2,3}K\\
        &=(x_1-x_2)^{2m+1}(x_1-x_3)^{2m+1}(Ks_{2,3}L-Ls_{2,3}K).
    \end{align*}
    So $As_{2,3}B-Bs_{2,3}A$ is zero if and only if $Ks_{2,3}L=Ls_{2,3}K$, or in other words if $Ks_{2,3}L$ is $s_{2,3}$-invariant. Let us assume that $Ks_{2,3}L$ is $s_{2,3}$-invariant. We will show that this implies $K=L$ up to a scalar, which contradicts $V$ and $W$ being distinct.
    
    Decompose $K$ into its irreducible factors, and let $Q$ be such a factor. We will show that $Q$ is also a factor of $L$. Since $Ks_{2,3}L$ is $s_{2,3}$-invariant, $s_{2,3}Q$ must be a factor of either $K$ or $s_{2,3}L$. If it is a factor of $s_{2,3}L$, then $P$ is a factor of $L$, as desired. Otherwise, $s_{2,3}Q$ is a factor of $K$. Then we also have that $s_{1,2}Q$ and $s_{1,2}s_{2,3}Q$ are factors of $K$ since $K$ is $s_{1,2}$-invariant. Then $s_{2,3}s_{1,2}Q$ is a factor of $Ks_{2,3}L$. If it is a factor in $s_{2,3}L$, then note that $s_{1,3}s_{2,3}s_{1,2}Q$ is a factor of $s_{2,3}L$ since $L$ being $s_{1,2}$-invariant implies $s_{2,3}L$ is $s_{1,3}$-invariant. But we have $s_{1,3}s_{2,3}s_{1,2}=s_{2,3}$, so $s_{2,3}Q$ is a factor of $s_{2,3}L$, which implies $Q$ is a factor of $L$, as desired. So we may assume $s_{2,3}s_{1,2}Q$ is a factor of $K$. Then we also have $s_{1,2}s_{2,3}s_{1,2}Q=s_{1,3}Q$ is a factor of $K$ since $K$ is $s_{1,2}$-invariant, and thus $K$ is divisible by $sQ$ for all $s\in S_3$. 
    
    So $K$ is divisible by the least common multiple of all of the $sQ$, which we will denote by $S$. $S_3$ fixes $S$ as a factor of $K$, so it must act on $S$ by scalars. It cannot act on $S$ trivially, as then $S$ would be a symmetric polynomial factor of $K$, and $K$ has no symmetric polynomial factors by Corollary \ref{cor:3.3}. So $S_3$ must act on $S$ by sign. Note also that we can only have one such factor of $K$ on which $S_3$ acts by sign, as otherwise their product would be a symmetric polynomial factor of $K$. So every factor of $K$ is also a factor of $L$ except for possibly a factor in the sign representation. We may apply everything above to $L$ as well to get the same result in the other direction, yielding a bijection between all irreducible factors of $K$ and $L$ except maybe a factor of the sign representation in either case. Since $K$ and $L$ are both $s_{1,2}$-invariant, acting on both of them by $s_{1,2}$ makes it clear that either they both have an extra factor of the sign representation or they both do not.
    
    If $K$ and $L$ do not have an additional factor in the sign representation, then we have a bijection between all irreducible factors of $K$ and $L$, and thus they are equal up to a scalar, as desired. Otherwise, note that by Proposition \ref{prop:2.2}, the sign factors lie in $Q_0(3,\k)$, and by Proposition \ref{prop:2.6} they are symmetric polynomial multiples of $(x_1-x_2)(x_2-x_3)(x_3-x_1)$. But these symmetric polynomial multiples must be 1, as otherwise $K$ and $L$ would be divisible by nonconstant symmetric polynomials, which violates Corollary \ref{cor:3.3}. So the sign factors are both $(x_1-x_2)(x_2-x_3)(x_3-x_1)$. Thus we still have a bijection between the factors of $K$ and $L$, so $K$ and $L$ are equal up to a scalar, as desired.
    
    Note that $As_{2,3}B-Bs_{2,3}A$ has degree $\deg V+\deg W$. Since $Q_m(3,\k)_\sgn$ is generated by a polynomial of degree $6m+3$ by Proposition \ref{prop:2.6}, we have $\deg V+\deg W\geq 6m+3$.
\end{proof}

The above lemma brings us closer to the proof of Theorem \ref{thm:3.1}, however we still need a couple more lemmas.

\begin{lemma}\label{lem:3.5}
    Assume that there exist generating representations $V,W$ of $Q_m(3,\k)_\std$ such that $\deg V+\deg W=6m+3$. Then $Q_m(3,\k)_\std$ is a free module over $\k[x_1,x_2,x_3]^{S_3}$ generated by $V$ and $W$.
\end{lemma}
\begin{proof}
    Assume for contradiction that there exists some other generating representation $U$ of \linebreak
    $Q_m(3,\k)_\std$. By Lemma \ref{lem:3.4} we have $\deg U\geq \deg W$. Let $A\in V_{1,2}^-,B\in W_{1,2}^-,C\in U_{1,2}^-$. Then note that by Lemma \ref{lem:3.4} we have
    $$As_{2,3}B-Bs_{2,3}A=c\prod_{i<j}(x_i-x_j)^{2m+1}$$ 
    where $c$ is a nonzero constant, since it is a degree $6m+3$ polynomial in $Q_m(3,\k)_\sgn$. We also have
    $$As_{2,3}C-Cs_{2,3}A=Q\prod_{i<j}(x_i-x_j)^{2m+1}$$
    for some symmetric polynomial $Q$. But then instead of $U$, we may take the representation generated by $cC-QB$ as the generating representation of $Q_m(3,\k)_\std$. But we have
    \begin{align*}
        As_{2,3}(cC-QB)-(cC-QB)s_{2,3}A&=c(As_{2,3}C-Cs_{2,3}A)-Q(As_{2,3}B-Bs_{2,3}A)\\
        &=(cQ-Qc)\prod_{i<j}(x_i-x_j)^{2m+1}\\
        &=0
    \end{align*}
    which contradicts Lemma \ref{lem:3.4}. If $Q_m(3,\k)_\std$ was not a free module then there would be nonzero symmetric polynomials $R_1,R_2$ such that $R_1A=R_2B$. But then we would have 
    $$R_1As_{2,3}(R_2B)-R_2Bs_{2,3}(R_1A)=R_1As_{2,3}(R_1A)-R_1As_{2,3}(R_1A)=0$$
    and 
    $$R_1As_{2,3}(R_2B)-R_2Bs_{2,3}(R_1A)=R_1R_2(As_{2,3}B-Bs_{2,3}A)$$
    which cannot happen since none of $R_1,R_2,as_{2,3}b-bs_{2,3}a$ are zero.
\end{proof}
\begin{thebibliography}{99}
\bibitem[2]{wangv2:berest2003cherednik} Berest Yu., Etingof P., Ginzburg V., Cherednik algebras and differential operators on quasi-invariants, Duke Mathematical Journal 118 (2003), 279-- 337.
\bibitem[6]{wangv2:etingof2002lectures} Etingof P., Strickland E., Lectures on quasi-invariants of Coxeter groups and the Cherednik algebra, arXiv preprint math/0204104 (2002).
\bibitem[8]{wangv2:felder2003action} Felder G., Veselov A., Action of Coxeter groups on m-harmonic polynomials and Knizhnik-Zamolodchikov equations, Moscow Math. J (2003).
\bibitem[10]{wangv2:humphreys_1990} Humphreys J.  E., Polynomial invariants of finite reflection groups, Reflection Groups and Coxeter Groups (1990), 49-- 86.
\bibitem[12]{wangv2:Ren_2020} Ren M., Xu X., Quasi-Invariants in Characteristic p and Twisted Quasi-Invariants, Symmetry, Integrability and Geometry: Methods and Applications (2020).
\end{thebibliography}
\endgroup
\begin{thebibliography}{99}
\bibitem[1]{Alperin_1986}
Alperin J.L., Local representation theory: {M}odular representations as an
 introduction to the local representation theory of finite groups,
 \textit{Cambridge Stud. Adv. Math.}, Vol.~11,
 \href{https://doi.org/10.1017/CBO9780511623592}{Cambridge University Press},
 Cambridge, 1986.

\bibitem[4]{calogero}
Calogero F., Solution of the one-dimensional {$N$}-body problems with quadratic
 and/or inversely quadratic pair potentials,
 \href{https://doi.org/10.1063/1.1665604}{\textit{J.~Math. Phys.}} \textbf{12}
 (1971), 419--436.

\bibitem[6]{cmp/1104179957}
Chalykh O.A., Veselov A.P., Commutative rings of partial differential operators
 and {L}ie algebras, \href{https://doi.org/10.1007/BF02125702}{\textit{Comm.
 Math. Phys.}} \textbf{126} (1990), 597--611.

\bibitem[8]{etingof2002lectures}
Etingof P., Strickland E., Lectures on quasi-invariants of {C}oxeter groups and
 the {C}herednik algebra, \textit{Enseign. Math.} \textbf{49} (2003), 35--65,
 \href{http://arxiv.org/abs/math.QA/0204104}{arXiv:math.QA/0204104}.

\bibitem[10]{MOSER1975197}
Moser J., Three integrable {H}amiltonian systems connected with isospectral
 deformations,
 \href{https://doi.org/10.1016/0001-8708(75)90151-6}{\textit{Adv. Math.}}
 \textbf{16} (1975), 197--220.

\bibitem[11]{Ren_2020}
Ren M., Xu X., Quasi-invariants in characteristic {$p$} and twisted
 quasi-invariants,
 \href{https://doi.org/10.3842/SIGMA.2020.107}{\textit{SIGMA}} \textbf{16}
 (2020), 107, 13~pages,
 \href{http://arxiv.org/abs/1907.13417}{arXiv:1907.13417}.

\bibitem[13]{wang2022explicit}
Wang F., Toward explicit {H}ilbert series of quasi-invariant polynomials in
 characteristic {$p$} and {$q$}-deformed quasi-invariants, \textit{New York~J.
 Math.} \textbf{29} (2023), 613--634,
 \href{http://arxiv.org/abs/2201.06111}{arXiv:2201.06111}.

\end{thebibliography}
\end{document}
